% appendix-numerical-error-tracking.tex — \input'd from appendix-numerical.tex
%
% How the numerical algorithm reports results: certified vs uncertain
% candidates, and when the answer is conclusive.
%
% Status: draft, not yet Jörn-reviewed.

After evaluating all candidate pairs $(S, \sigma)$,
the numerical algorithm must report a capacity value with
appropriate certainty qualifiers.
Near-singular systems have already been resolved by
Algorithm~\ref{alg:near-singular-handling}: each pair is
either dismissed (value covered by a smaller pair) or its
computed~$\beta_0$ is used directly.
The only remaining source of uncertainty is the
three-valued predicates of
Section~\ref{app:three-valued}.

\paragraph{Candidate classification.}
For each non-dismissed pair $(S, \sigma)$, the algorithm
stores the full candidate $(S, \sigma, \hat\beta, \hat Q)$:
the pair, the computed coefficient vector, and the
objective value $\hat Q = \hat\beta^\top H\,\hat\beta$.
Each candidate is classified as:
\begin{description}
\item[Certified.]
  All predicates returned $\textsc{true}$.
  The candidate is genuinely admissible and~$\hat Q$ is
  the value of a feasible critical point.
\item[Uncertain.]
  Some predicate returned $\textsc{unknown}$
  (typically a $\hat\beta_i$ near zero).
  The pair may or may not satisfy the exact constraints,
  so~$\hat Q$ may come from an inadmissible point.
  The stored~$\hat\beta$ enables downstream investigation
  (e.g.\ more expensive tests to resolve the
  $\textsc{unknown}$ verdicts).
\end{description}

\paragraph{Final answer.}
Following the maximum tracking principle
(Section~\ref{app:three-valued}):
let $(S_{\text{cert}}, \sigma_{\text{cert}},
\hat\beta_{\text{cert}}, Q_{\text{cert}})$ be the
certified candidate with the largest~$Q$, and let
$\mathcal{U} = \{(S, \sigma, \hat\beta, \hat Q) :
\text{uncertain with } \hat Q > Q_{\text{cert}}\}$
be the uncertain candidates above the certified value.
If no certified candidate exists, the algorithm reports
only the uncertain candidates and flags the result for
investigation.
Otherwise:
\begin{enumerate}
\item
  $\mathcal{U} = \emptyset$:
  the certified candidate dominates all uncertain ones.
  The capacity is
  $c_{\mathrm{EHZ}}(K) = 1 / (2\, Q_{\text{cert}})$
  with high confidence.
\item
  $\mathcal{U} \neq \emptyset$:
  some uncertain candidate has a higher~$Q$ value.
  The certified value gives a reliable lower bound on the
  maximum~$Q$ (upper bound on capacity):
  $c_{\mathrm{EHZ}}(K) \leq 1 / (2\, Q_{\text{cert}})$.
  If any uncertain candidate in~$\mathcal{U}$ is truly
  admissible, the true capacity is lower.
  The algorithm returns both $Q_{\text{cert}}$ and
  $\mathcal{U}$ for downstream investigation.
\end{enumerate}

\begin{remark}[Empirical error tracking]\label{rem:empirical-errors}
In practice, we track error bounds empirically in the Rust
implementation rather than deriving tight analytic bounds
for every operation.
The SVD decomposition accounts for roughly 80\% of the
computation time; the error tracking happens \emph{after}
the SVD and does not require modifying the SVD computation
itself.
The implementation tracks the best certified and best
uncertain candidates separately, evaluating the dismissal
error bound (Remark~\ref{lem:dismissal-error-bound})
for each dismissed system.
This makes the error tracking essentially free in terms of
additional computational cost.
\end{remark}
